\section{Introduction and welcome}

The workshop, organized by the COMETA COST Action (Comprehensive Multiboson
Experiment-Theory Action, CA22130), took place at the Aula Magna of the Agraria
Department of the Universit\`a di Perugia. In the opening session, Valentina
Mariani (Universit\`a e INFN, Perugia) welcomed the participants and presented
the practical organization and the goals of the meeting: to bring together the
theoretical, experimental, and machine-learning communities working on boosted
hadronic decays of electroweak bosons, with a spin on polarization tagging and
on vector-boson-fusion/scattering topologies.

The physics case underlying the whole workshop can be stated concisely.
Longitudinally polarized $W$ and $Z$ bosons are the physical remnants of the
Higgs mechanism; in the Standard Model (SM) the scattering of longitudinal
bosons ($V_L V_L \to V_L V_L$) is unitarized by delicate gauge and Higgs
cancellations, so measuring polarized (in particular doubly longitudinal)
diboson production and scattering is a direct probe of electroweak symmetry
breaking (EWSB) and a sensitive portal to new physics. All measurements to date
use fully leptonic decays, which are clean but statistically starved: the
longitudinal fraction in VBS is only $\mathcal{O}(5\text{--}10)\%$ of an already
rare process. Hadronic boson decays offer an order-of-magnitude larger
branching fraction and full kinematic reconstruction in the boosted regime,
where the decay is contained in a single large-radius jet, but require new
tools---polarization-aware jet taggers, decorrelation techniques, precise Monte
Carlo modeling of polarized production and hadronic decays---which formed the
core of the two-day program.
